\customchapter{ABSTRACT}

\begin{center}
\large \vspace{-1.5cm} \textbf{\MakeUppercase{DARL: Drift-Aware Reinforcement Learning for Selective Updating of Two-Stage Tabular ML Pipelines}}
\end{center}

Machine learning models deployed in production lose quality when incoming data no longer
resembles the data they were trained on. In a two-stage pipeline, made of a data
preparation stage and a predictive model, this degradation may originate in either stage.
Yet the usual response is to retrain the whole pipeline, which is costly and not always
necessary. This thesis proposes DARL (Drift-Aware Reinforcement Learner), an approach that
treats pipeline maintenance as a sequence of decisions. Each day, a reinforcement learning
agent reviews monitoring signals and chooses between doing nothing, correcting only the data
preparation stage, retraining only the model, or retraining the whole pipeline. Since the
agent cannot directly observe the type of shift taking place, the problem is formulated as
a Partially Observable Markov Decision Process and solved with a Deep Q-Network (DQN). The
reward combines the improvement in the area under the precision-recall curve (AUPRC) with
the cost of each action, and is computed with a one-day delay, once the true labels become
available. The case study is early sepsis prediction using data from the
PhysioNet/Computing in Cardiology Challenge 2019: a pipeline trained on patients from one
hospital is used in a simulated 200-bed intensive care unit of another hospital, where
controlled shifts in vital sign measurement, in sepsis labeling, and their combination are
injected. The learned policy is contrasted with reference rules such as never acting,
always retraining, acting at random, or following a fixed threshold.

\noindent \textbf{Keywords:}\\
\noindent Distribution shift; Reinforcement learning; Model maintenance; Sepsis prediction
